\begin{table}[htbp]
\centering
\caption{Survey-weighted logistic regression predicting usable child height data (pooled DHS surveys, 2011--2022)}
\label{tab:logit_main}
\begin{threeparttable}


\small
\renewcommand{\arraystretch}{1.12}
\setlength{\tabcolsep}{6pt}
\begin{adjustbox}{max width=\textwidth}
\begin{tabularx}{\textwidth}{X l r r r}
\toprule
\textbf{Variable} & \textbf{Level} & \textbf{Adj. OR} & \textbf{95\% CI} & \textbf{p-value} \\
\midrule

\addlinespace[3pt]
\multicolumn{5}{l}{\textbf{Child characteristics}}\\
\addlinespace[2pt]

Child age (months) & Continuous & 1.012 & (1.011, 1.013) & \(<0.001\) \\

Sex of child & Male   & 1.00 &  &  \\
             & Female & 1.079 & (1.046, 1.115) & \(<0.001\) \\

Ill in past 2 weeks & No  & 1.00 &  &  \\
                   & Yes & 1.340 & (1.287, 1.395) & \(<0.001\) \\

Number of children in household & Continuous & 0.966 & (0.948, 0.985) & 0.001 \\

\addlinespace[4pt]
\multicolumn{5}{l}{\textbf{Maternal characteristics}}\\
\addlinespace[2pt]

Mother's age (years) & Continuous & 0.999 & (0.995, 1.003) & 0.675 \\

Mother's education & No education & 1.00 &  &  \\
                  & Primary   & 1.162 & (1.098, 1.229) & \(<0.001\) \\
                  & Secondary & 1.164 & (1.111, 1.219) & \(<0.001\) \\
                  & Higher    & 1.011 & (0.944, 1.083) & 0.759 \\

Currently pregnant & No  & 1.00 &  &  \\
                  & Yes & 1.169 & (1.092, 1.250) & \(<0.001\) \\

\addlinespace[4pt]
\multicolumn{5}{l}{\textbf{Household characteristics}}\\
\addlinespace[2pt]

Wealth index & Poorest & 1.00 &  &  \\
            & Poorer  & 1.052 & (1.003, 1.104) & 0.038 \\
            & Middle  & 1.005 & (0.952, 1.061) & 0.858 \\
            & Richer  & 0.979 & (0.915, 1.049) & 0.547 \\
            & Richest & 0.776 & (0.726, 0.830) & \(<0.001\) \\

Household head & Male   & 1.00 &  &  \\
               & Female & 1.039 & (0.989, 1.092) & 0.133 \\

Residence & Urban & 1.00 &  &  \\
          & Rural & 1.344 & (1.274, 1.417) & \(<0.001\) \\

\bottomrule
\end{tabularx}
\end{adjustbox} 

\begin{tablenotes}[flushleft]
\footnotesize
\item \textbf{Model information.} Observations: 509{,}006. PSUs (clusters): 62{,}643. Outcome: usable height data (1 = usable; 0 = missing/implausible). Odds ratios are derived from survey-weighted logistic regression models accounting for DHS sampling weights, clustering at the PSU level, and stratification.
\end{tablenotes}

\end{threeparttable}
\end{table}